\documentclass[a4paper,11pt]{article}
% Shared preamble — % Shared preamble — % Shared preamble — \input{../shared/preamble} in every document
% Each document must declare \documentclass[a4paper,11pt]{article} before \input-ing this file.
\usepackage[T1]{fontenc}
\usepackage[utf8]{inputenc}
\usepackage[polish,english]{babel}
\usepackage[top=2cm,bottom=2cm,left=2cm,right=2cm]{geometry}
\usepackage{booktabs}
\usepackage{array}
\usepackage{tabularx}
\usepackage{longtable}
\usepackage{multirow}
\usepackage{hhline}
\usepackage{enumitem}
\usepackage{xcolor}
\usepackage{fancyhdr}
\usepackage{graphicx}
\usepackage{lastpage}

\definecolor{pzzblue}{RGB}{0,70,127}

\pagestyle{fancy}
\fancyhf{}
\renewcommand{\headrulewidth}{0.4pt}
\fancyfoot[C]{\thepage/\pageref{LastPage}}

% Helper: underlined fill-in field of given width
\newcommand{\field}[1]{\underline{\hspace{#1}}}
\newcommand{\fieldline}{\noindent\rule{\linewidth}{0.4pt}}
 in every document
% Each document must declare \documentclass[a4paper,11pt]{article} before \input-ing this file.
\usepackage[T1]{fontenc}
\usepackage[utf8]{inputenc}
\usepackage[polish,english]{babel}
\usepackage[top=2cm,bottom=2cm,left=2cm,right=2cm]{geometry}
\usepackage{booktabs}
\usepackage{array}
\usepackage{tabularx}
\usepackage{longtable}
\usepackage{multirow}
\usepackage{hhline}
\usepackage{enumitem}
\usepackage{xcolor}
\usepackage{fancyhdr}
\usepackage{graphicx}
\usepackage{lastpage}

\definecolor{pzzblue}{RGB}{0,70,127}

\pagestyle{fancy}
\fancyhf{}
\renewcommand{\headrulewidth}{0.4pt}
\fancyfoot[C]{\thepage/\pageref{LastPage}}

% Helper: underlined fill-in field of given width
\newcommand{\field}[1]{\underline{\hspace{#1}}}
\newcommand{\fieldline}{\noindent\rule{\linewidth}{0.4pt}}
 in every document
% Each document must declare \documentclass[a4paper,11pt]{article} before \input-ing this file.
\usepackage[T1]{fontenc}
\usepackage[utf8]{inputenc}
\usepackage[polish,english]{babel}
\usepackage[top=2cm,bottom=2cm,left=2cm,right=2cm]{geometry}
\usepackage{booktabs}
\usepackage{array}
\usepackage{tabularx}
\usepackage{longtable}
\usepackage{multirow}
\usepackage{hhline}
\usepackage{enumitem}
\usepackage{xcolor}
\usepackage{fancyhdr}
\usepackage{graphicx}
\usepackage{lastpage}

\definecolor{pzzblue}{RGB}{0,70,127}

\pagestyle{fancy}
\fancyhf{}
\renewcommand{\headrulewidth}{0.4pt}
\fancyfoot[C]{\thepage/\pageref{LastPage}}

% Helper: underlined fill-in field of given width
\newcommand{\field}[1]{\underline{\hspace{#1}}}
\newcommand{\fieldline}{\noindent\rule{\linewidth}{0.4pt}}


\fancyhead[L]{\textbf{KARTA REJSU} — Morze Bałtyckie}
\fancyhead[R]{PZŻ}

\begin{document}
\selectlanguage{polish}

\begin{center}
    {\Large\bfseries\color{pzzblue} KARTA REJSU}\\[4pt]
    {\large Polski Związek Żeglarski}\\[2pt]
    {\normalsize Strefa: \textbf{Morze Bałtyckie}}
\end{center}

\vspace{0.5cm}
\noindent\rule{\linewidth}{1pt}
\vspace{0.3cm}

% --- Dane jachtu ---
\subsection*{Dane jachtu}
\begin{tabularx}{\linewidth}{|l|X|l|X|}
\hline
Typ jachtu & \field{4cm} & Nazwa jachtu & \field{4cm} \\[6pt]
\hline
Nr rejestracyjny & \field{4cm} & Bandera & \field{4cm} \\[6pt]
\hline
Armator / Czarterujący & \multicolumn{3}{X|}{\field{10cm}} \\[6pt]
\hline
Długość [m] & \field{2.5cm} & Powierzchnia żagli [m²] & \field{2.5cm} \\[6pt]
\hline
Silnik pomocniczy & \field{3cm} & Moc [KM/kW] & \field{3cm} \\[6pt]
\hline
\end{tabularx}

\vspace{0.4cm}

% --- Dane kapitana ---
\subsection*{Kapitan / Skipper}
\begin{tabularx}{\linewidth}{|l|X|l|X|}
\hline
Nazwisko i imię & \field{5cm} & PESEL / nr dok. & \field{4cm} \\[6pt]
\hline
Nr patentu / uprawnień & \field{4cm} & Stopień & \field{4cm} \\[6pt]
\hline
Adres & \multicolumn{3}{X|}{\field{10cm}} \\[6pt]
\hline
Telefon & \field{4cm} & E-mail & \field{5cm} \\[6pt]
\hline
\end{tabularx}

\vspace{0.4cm}

% --- Trasa ---
\subsection*{Trasa rejsu}
\begin{tabularx}{\linewidth}{|l|X|l|X|}
\hline
Port wyjścia & \field{4cm} & Data wyjścia & \field{3cm} \\[6pt]
\hline
Port docelowy & \field{4cm} & Data powrotu & \field{3cm} \\[6pt]
\hline
Planowane porty pośrednie & \multicolumn{3}{X|}{\field{10cm}} \\[6pt]
\hline
Liczba mil morskich & \field{3cm} & Liczba dób na morzu & \field{3cm} \\[6pt]
\hline
\end{tabularx}

\vspace{0.4cm}

% --- Informacje strefowe Bałtyk ---
\subsection*{Informacje strefowe — Bałtyk}
\begin{itemize}[noitemsep]
    \item Akwen: \field{5cm} (np.\ Zatoka Gdańska, Zalew Wiślany, Bałtyk otwarty)
    \item Kapitanat portu wyjścia: \field{6cm}
    \item Nr VHF dyżurne (BSH/MRCC): Ch~16 / Ch~\field{1.5cm}
    \item Zasięg pływania: \field{3cm} Mm od brzegu
    \item Prognoza IMGW (data/godzina): \field{4cm}
\end{itemize}

\vspace{0.4cm}

% --- Załoga ---
\subsection*{Załoga}
\begin{tabularx}{\linewidth}{|c|X|l|l|l|}
\hline
\textbf{Lp.} & \textbf{Nazwisko i imię} & \textbf{Data ur.} & \textbf{Nr dok.} & \textbf{Funkcja} \\
\hline
1 & & & & \\[5pt]
\hline
2 & & & & \\[5pt]
\hline
3 & & & & \\[5pt]
\hline
4 & & & & \\[5pt]
\hline
5 & & & & \\[5pt]
\hline
6 & & & & \\[5pt]
\hline
7 & & & & \\[5pt]
\hline
8 & & & & \\[5pt]
\hline
\end{tabularx}

\vspace{0.4cm}

% --- Wyposażenie bezpieczeństwa ---
\subsection*{Wyposażenie bezpieczeństwa (Bałtyk)}
\begin{tabularx}{\linewidth}{|X|c|X|c|}
\hline
\textbf{Element} & \textbf{Tak/Nie} & \textbf{Element} & \textbf{Tak/Nie} \\
\hline
Kamizelki ratunkowe (szt.) & \field{1cm} & Tratwa ratunkowa & \\[4pt]
\hline
Koło ratunkowe & & Lina asekuracyjna / uprząż & \\[4pt]
\hline
Rakiety sygnałowe & & EPIRB / PLB & \\[4pt]
\hline
Radio VHF & & Gaśnica & \\[4pt]
\hline
Apteczka & & Instrukcja bezpieczeństwa & \\[4pt]
\hline
\end{tabularx}

\vspace{1cm}

% --- Podpisy ---
\noindent
\begin{tabularx}{\linewidth}{X X}
Podpis kapitana: & Podpis armatora / przedstawiciela: \\[1.2cm]
\rule{6cm}{0.4pt} & \rule{6cm}{0.4pt} \\
Data: \field{3cm} & Data: \field{3cm}
\end{tabularx}

\end{document}
